\subsection{Physics modelling}
\label{subsection_physics_modelling}

The key idea is to represent physical entities using objects that are represented directly
in the underlying algebra. They have a geometric interpretation and can be directly
manipulated by applying algebraic operations. \\

$EGA$ is limited to objects that contain the origin. Thus, it can be used to directly
model vectors (i.e. directions), lines, and planes through the orgin. In addition, it can
model complex numbers, quaternions and thus rotations in $2D$ and $3D$ well. To realize
them, it uses sandwich products based on the geometric product for calculating the
transformation resulting from two consequtive reflections either across lines ($2d$-case)
or across planes ($3D$-case). The resulting transformations are rotations. The advantage
of using geometric algebra to model those rotations is the intuitive modelling, due to
geometric meaning of the objects and the fact that the sandwich products can be used to
transform different objects, like vectors, bivectors or rotation operators using the same
sandwich product formulas in a unified fashion regardless of the type of object to be
transformed.\\

In addition, $PGA$ is capable to model objects that do not necessarily contain the origin.
Thus it can be used to directly model points, directions, lines, and planes that may or
may not contain the origin. Translations can be implemented in $PGA$ directly as
orthogonal transformation using a sandwich product. The operation is realized by two
consequtive reflections across parallel lines ($2D$) or planes ($3D$). It is based on the
regressive form of the geometric product. The reason for this deviation from $EGA$, which
uses the geometric product for that purpose, is that the Euclidean isometries (rotation,
translation, and screw motion) can only be performed in a space that is dual to the space
we use for modelling our projective objects (for explanation
see~\cite{Lengyel_pga-illuminated:2024}). The use of the regressive geometric product on
these objects is equivalent to first transforming the object into dual space by applying
the complement operation, then using the geometric product in that space for the intended
transformation and finally transforming the resulting transformed object back into the
original space, by applying the inverse complement operation. Thus the sandwich product
transformation formulas in $PGA$ differ from the formulas used in $EGA$ by using the
regressive geometric product instead of the geometric product, but are otherwise the same.
\\

The additional degree of freedom in $PGA$ is possible due to embedding the Euclidean space
of dimension $n$ into a projective space of dimension $n+1$, e.g. two-dimensional
Euclidean space is modelled by embedding it in a three-dimensional projective space, and
three-dimensional Euclidean space is modelled in a four-dimensional projective space. In
the $n+1$-dimensional projective space all algebraic objects contain the origin (exactly
as in $EGA$), but in the modelled $n$-dimensional Euclidean space the objects do not need
to contain the origin. They can be located anywhere in the modelled space including
postions at infinity. In projective geometry the objects must be projected into the
modelled space. Thus the objects of the embedding space have one dimension more than the
corresponding objects in modelled space. A projective point (which models a scalar or
$0d$-object) is a vector representing a ray in the embedding space (a $1d$-object), and (a
$1d$-object) like a line in modelled space is created by a bivector in the embedding
projective space (a $2d$-object), and so on for higher grade objects. The projection
reduces the dimension of the object in the embedding space by one. \\

Due to its extended expressiveness $PGA$ is suitable to model physics of mass points and
rigid bodies. For example it can model position, velocity, acceleration, force, torque,
momentum, angular momentum, etc., and thus it can be used to describe kinematics and
dynamics of mass points and rigid bodies. The approach used here is inspired by
\cite{Plane-based_PGA_Dorst-DeKennik:2022} and
\cite{Dynamics_in_PGA_plane-based_Dorst-DeKennik:2023}, but it does not use the
plane-based ansatz with planes as vectors, which are based on dual representations of the
objects to be modelled. It rather uses the approach as described by
\cite{Lengyel_pga-illuminated:2024} using projective geometric algebra based on direct
representations of points, lines and planes, who also gives reasons for choosing the
latter approach and explains the background in more detail compared to the short
indications given here. \\

Objects modelled in projective space can be multiplied by a scalar value without changing
the \emph{geometric} meaninig of the object in the modelled Euclidean space, i.e. its
interpretation after the projection. Equality of modelled geometric objects is thus often
only defined up to a scalar factor. This kind of equality up to a scalar factor is defined
as congruence. Therefore, it is important for modelling of physical facts that unitization
is used consciously to avoid false interpretations. For example, a heavy point can be
modelled by multiplying a \emph{unitized} point with the scalar value of its mass, but of
course the multiplication with $m$ must happen after unitization in order to keep the
scalar factor $m$ in the calculation. Other geometric objects need to be unitized in order
to focus on the intended geometric operation and avoid additional scaling, i.e. the
rotation operators. They must be unitized before use as well. Thus it is advisable to
think through the geometric operations and their physical interpretation well to avoid
suprises and to be consistent.

\cite{Lengyel_poor-foundations_GA:2024} points out potential mathematical inconsistencies
in applications of geometric algebra as often used up-to-now. Further hints on
inconsistencies in some GA approaches can be found as well in
\cite{Kritchevsky_case_against_GA:2024}. The goal of the this work is to be as consistent
and mathematically rigoros as possible implementing GA functionality, and at the same time
providing an intuitively applicable framework for modelling of physics and applications.
\\


\subsubsection{Modelling force and torque}
\label{subsubsection_modelling_force_and_torque}

\cite{Browne_Grassmann-Algebra_Vol1:2012} (see sections 4.1-4.5) is a very good source
showing how to use Grassmann algebra and projective geometry in order to model forces and
moments. Classical approaches based-on vector algebra only model forces by providing
vectors of adequate magnitude and direction to represent the force. However, the
physically very meaningful ``line of action'' is not modelled mathematically at all in the
classical approach. It rather needs to be derived from the description of the context in
order to know where the force is actually acting. \\

This changes in $PGA$. Any force can be directly modelled by its line of action which
encodes the force \emph{and} the torque generated by the force at the origin within one
object. A force $f$ is modelled by a line $F$. If $f$ is the force vector of corresponding
magnitude and direction, and $Q$ is any unitized projective point on its line of action
then $F$ is the resulting line of action:
\begin{equation}
F = Q \wedge f
\label{eq:line_of_action}
\end{equation}
The resulting bivector $F$ represents a line, a geometric object which contains all
physically relevant information of the force $f$ and the torque generated by $f$. For that
reason~\cite{Dynamics_in_PGA_plane-based_Dorst-DeKennik:2023} call this object ``forque''.
It can directly be used to recover the force vector and to calculate the resulting free
torque by the force either at the origin or at any point in space as shown below. \\

But first let's have a closer look how this is accomplished. We can separate $F$ into its
components. Each unitized projective point $Q$ is the sum of the origin $O$ and a vector
$q$ parallel to the projective plane pointing from $O$ to $Q$, such that $Q = O + q$
(``projective point split''). Thus $F$ can be split into two bivector components
\begin{equation}
    F = Q \wedge f = (O + q) \wedge f =
    \underbrace{O \wedge f}_{\text{=force}} + \underbrace{q \wedge f}_{=torque}
    = \weight{F} + \bulk{F}
\label{eq:line_of_action_components}
\end{equation}
where the first term encodes the force as within the weight components of the bivector,
while the second term encodes the resulting free torque relative to the origin caused be
the action of $f$ at $Q$ -- or alternatively anywhere along the line of action -- within
the bulk components of the bivector. \\

To see this even more clearly, let's look at the component equations of both terms written
as resulting bivectors for the case of $2D$ $PGA$ (the corresponding basis bivector
compontents are shown on the right-hand side in square brackets for reference):
\begin{subeqnarray}
    O \wedge f & = &
    \begin{pmatrix}
        0 \\ 0 \\ 1 
    \end{pmatrix}
    \wedge
    \begin{pmatrix}
         f_x \\ f_y \\ 0 
    \end{pmatrix}
    =
    \begin{pmatrix}
         -f_y \\ f_x \\ 0 
    \end{pmatrix}
    \qquad \qquad \quad \;\;
    \begin{bmatrix}
         \bv{23} \\ \bv{31} \\ \bv{12} 
    \end{bmatrix}
    \slabel{eq:pga2dp_force_line_force_component} \\[10 pt]
        q \wedge f & = &
    \begin{pmatrix}
        q_x \\ q_y \\ 0 
    \end{pmatrix}
    \wedge
    \begin{pmatrix}
         f_x \\ f_y \\ 0 
    \end{pmatrix}
    =
    \begin{pmatrix}
         0 \\ 0 \\ q_x f_y - q_y f_x
    \end{pmatrix}
    \qquad
    \begin{bmatrix}
         \bv{23} \\ \bv{31} \\ \bv{12} 
    \end{bmatrix}
    \slabel{eq:pga2dp_force_line_torque_component}
\end{subeqnarray}
where equation (\ref{eq:pga2dp_force_line_force_component}) shows how the force is encoded
in the first two components of the bivector and  equation
(\ref{eq:pga2dp_force_line_torque_component}) shows how the torque is encoded in the third
component of the bivector $F$. \\

For $3D$ $PGA$ the encoding is similar, but of course now requires six bivector components
instead of three (the corresponding basis bivector compontents are shown on the right-hand
side in square brackets for reference):
\begin{subeqnarray}
    O \wedge f & = &
    \begin{pmatrix}
        0 \\ 0 \\ 0 \\ 1 
    \end{pmatrix}
    \wedge
    \begin{pmatrix}
         f_x \\ f_y \\ f_z \\ 0 
    \end{pmatrix}
    =
    \begin{pmatrix}
         f_x \\ f_y \\ f_z \\ 0 \\ 0 \\ 0
    \end{pmatrix}
    \qquad \qquad \quad \quad \, %\;\;
    \begin{bmatrix}
         \bv{41} \\ \bv{42} \\ \bv{43} \\ \bv{23} \\ \bv{31} \\ \bv{12} 
    \end{bmatrix}
    \slabel{eq:pga3dp_force_line_force_component} \\[10 pt]
        q \wedge f & = &
    \begin{pmatrix}
        q_x \\ q_y \\ q_z \\ 0 
    \end{pmatrix}
    \wedge
    \begin{pmatrix}
         f_x \\ f_y \\ f_z \\ 0 
    \end{pmatrix}
    =
    \begin{pmatrix}
         0 \\ 0 \\ 0 \\ q_y f_z - q_z f_y \\ q_z f_x - q_x f_z \\ q_x f_y - q_y f_x
    \end{pmatrix}
    \qquad
    \begin{bmatrix}
         \bv{41} \\ \bv{42} \\ \bv{43} \\ \bv{23} \\ \bv{31} \\ \bv{12} 
    \end{bmatrix}
    \slabel{eq:pga3dp_force_line_torque_component}
\end{subeqnarray}
 \\

The force vector $f$ itself can be retrieved from $F$ by using the attitude function
$att()$:
\begin{equation}
    f = att(F) = F \vee \overline{\bv{n}}
    \label{eq:force_attitude}
\end{equation}
with $\bv{n}$ denoting the basis vector of the projective dimension that squares to zero
in the respective algebra. The torque generated relative to the origin $O$ is encoded in
the bulk of the bivector:
\begin{equation}
    M_{O} = bulk(F) = \bulk{F}
    \label{eq:force_torque_origin}
\end{equation}
while the resulting torque from the application of $f$ at any given point $R$ in space can
be retrieved from $F$, either using any known point $Q$ on $F$\footnote{Which is always
available via the function $support(F)$. It provides the point on $F$ with the shortest
distance to the origin $O$ by orthogonally projecting the origin onto $F$.}, or
alternatively by projecting $R$ onto $F$ and using the perpendicular distance vector
$d_\perp$ from $R$ towards the line $F$ in order to calculate the resulting free torque
$M_F(R)$ as
\begin{subeqnarray}
    M_F(R) & = &  \bulk{\left[ (Q-R) \wedge f \right]}
             = \bulk{\left[ (Q-R) \wedge att(F) \right]}
             \slabel{eq:action_line_moment_simple} \\
    M_F(R) & = & \bulk{\left[ d_\perp \wedge f \right]}
             = \bulk{\left[ (F \vee (R \wedge \rightweightdual{F}) - R)
             \wedge att(F) \right]}
    \slabel{eq:action_line_moment_projection}
\end{subeqnarray}
where equation~(\ref{eq:action_line_moment_simple}) is the formula with the least
computational effort using a known point $Q$ on $F$, while
equation~(\ref{eq:action_line_moment_projection}) uses the othogonal projection of $R$
onto $F$ to calculate the perpendicular distance vector $d_\perp$. \\

The resulting total force and torque that act on an object are then given by the sum
\begin{equation}
    F_{total} = \sum_i{F_i}
    \label{eq:total_sum_forque}
\end{equation}
This resulting sum will be used later on for calculation of dynamics. \\

Please remember:
\begin{itemize}
    \item Neither a vector nor a bivector are located anywhere. They can be moved around
    freely. A line however represents a bound vector that is bound by its line. Thus the
    vector can only be moved along its line without changing the geometric and physical
    meaning (see~\cite{Browne_Grassmann-Algebra_Vol1:2012}, p. 169).
    \item A free torque is modelled by a bivector with force components being zero, i.e.
    the force vector has no direction and zero length. But the bivector still will contain
    a resulting torque value.
    \item The torque calculated from either equation (\ref{eq:action_line_moment_simple})
    or (\ref{eq:action_line_moment_projection}) is necessarily a free torque as can be
    seen from the formulas. The origin does not occur in the result anymore due to the
    difference of the two projective points involved.
    \item The sum of forces and their induced torque values can be directly created by the
    sum of their respective lines of action, i.e. by summing the bivectors representing
    those lines. The resulting lines of action contain all relevant information incl. the
    resulting vector sum of the forces and the resulting torque values from the summation.
    There is no need to calculate these separately. The result will always be a resulting
    bivector that potentially contains a bound vector and a resulting torque.
    (see~\cite{Browne_Grassmann-Algebra_Vol1:2012}, p. 176).
\end{itemize}

% The geometrical situation is depicted in figure XXX for a two-dimensional case.


\subsubsection{Modelling linear and angular momentum}
\label{subsubsection_modelling_momentum}

Momentum in the classical approach is split into linear momentum $\vec{p} = m\vec{v}$ and
angular momentum $\vec{L} = \vec{x} \times \vec{p} = \vec{x} \times m \vec{v}$. Both are
represented as vectors, but they represent entities with different physical meaning. So
they cannot be combined in one object in the classical approach. \\

The $PGA$ approach to momentum is to define both the linear and the angular aspects within
one entity. This is in complete analogy to what we have seen in the last subsection for
force and torque resulting in a force line of action. \\

Projective momentum $P$ of a point mass $m$ at position $X$ that moves with velocity
$\dot{X}$ is defined as a bivector $P$ where
\begin{equation}
    P = X \wedge m \dot{X}
    \label{momentum_definition}
\end{equation}
It defines the overall momentum at position $X$ including linear and angular momentum in
the bivector $P(X)$, which can be shown by splitting the projective point $X$ into its
components $X = O + x$ (``projective point split''). In the $3D$-case with the origin $O =
O_{3dp} = (0,0,0,1)$ and the direction vector $x$, where $x = (X_x, X_y, X_z, 0)$, i.e.
specifically with $X_w = 0$, which represents a direction towards a point at infinity,
this results in
\begin{subeqnarray}
    P(X) &=& X \wedge m \dot{X} = (O + x) \wedge m \dot{X} \nonumber \\
         &=& O \wedge m \dot{X} + x \wedge m \dot{X} \nonumber \\
         &=& O \wedge m \dot{(O + x)} + x \wedge m \dot{(O + x)} \nonumber \\
         &=& O \wedge m \dot{x} + x \wedge m \dot{x} \\
         &=& \underbrace{\quad\weight{P}\quad}_{\text{linear}}
             \;+\underbrace{\quad\bulk{P}\quad}_{\text{angular}}
\end{subeqnarray}
$P$ defines an instantaneous momentum line where the weight part $\weight{P}$ of the
bivector $P$ encodes the linear momentum and the bulk part $\bulk{P}$ encodes the angular
momentum with respect to the origin.\footnote{The ``projective point split'' works in this
way, since the origin is a destinguishable point. This is in contrast to $EGA$ or the
classical approach where the orgin is a null-vector, which has no lasting effect if it is
added to another point. This is one reason why $PGA$ can combine linear and angular
aspects as separate physical effects in a single algebraic object. The other is that the
required 6 degrees of freedom can be encoded in a bivector.} \\

The physical aspects encoded in the momentum line $P$ can be retrieved in complete analogy
to the case of force and torque in subsection
\ref{subsubsection_modelling_force_and_torque}: \\

The linear momentum vector $p$ itself can be retrieved from $P$ by using the attitude
function $att()$:
\begin{equation}
    p = m v = m \dot{x} = att(P) = P \vee \overline{\bv{n}}
    \label{eq:linear_momentum_from_attitude}
\end{equation}
with $\bv{n}$ denoting the basis vector of the projective dimension that squares to zero
in the respective algebra. The angular momentum $L$ generated relative to the origin $O$
is encoded in the bulk of the bivector:
\begin{equation}
    L_{O} = x \wedge m \dot{x} = bulk(P) = \bulk{P}
    \label{eq:angular_momentum_wrt_origin}
\end{equation}
while the resulting angular momentum relative to any given point $R$ in space can be
retrieved from the momentum line $P$ by
\begin{equation}
    L(R) = \bulk{\left[ (X-R) \wedge m\dot{X} \right]}
         = \bulk{\left[ (X-R) \wedge att(P) \right]}
           \slabel{eq:action_momentum_line_simple}
\end{equation}
using $X$ as a known point on $P$. This provides the resulting angular momentum $L(R)$
with respect to $R$. The momentum line encodes the resulting momentum and its linear and
angular aspects relative to any arbitrarily chosen point. \\

The equations (\ref{eq:linear_momentum_from_attitude}) and
(\ref{eq:angular_momentum_wrt_origin}) given above permit to separate the linear and
angular effects as in the classical approach. However, these aspects need not be separated
at all and can be calculated together in $PGA$.\\

Using eqn.~(\ref{momentum_definition}) and the definition $X=O + x$ with $X$ as a unitized
projective point, $O$ as the chosen origin and $x$ as the direction vector pointing from
$O$ to $X$ we can write
\begin{equation}
    P_O(X) = X \wedge m \dot{X} = (O + x) \wedge m \dot{(O + x)}
    \label{momentum_written_with_O}
\end{equation}
Compared to eqn.~(\ref{momentum_definition}) an index $O$ was added to $P$, since the
formula (\ref{momentum_written_with_O}) seems to confirm that the momentum depends on the
choice of the origin. Now it is important to demonstrate  that this is \emph{not} the
case. We are going to prove that the described physics does not depend on an arbitrarily
chosen origin and not on the resulting coordinates used for its description. \\

\emph{First approach}: Use a transformation of the standpoint from which we describe the
physical situation without actually moving the coordinate system (we still want use the
chosen coordinate system to describe the vectors and bivectors used): We go to a virtual
origin $O' = Y$, but describe the situation in the same coordinate system. This results in
describing the origin as $O = Y - y$ as seen from $Y$. Inserting that definition into
eqn.~(\ref{momentum_written_with_O}) and using the geometric vector identity
\begin{equation}
X = Y + (X-Y) = Y + (O+x-(O+y)) = Y + (x-y)
\label{eq:identity_transformation_X_Y}
\end{equation}
we express $P(X)$ as viewed from $Y$ by reformulating the equation for $P$ relative to the
origin $O$
\begin{subeqnarray}
    P_O(X) &=& X \wedge m \dot{X} = (O + x) \wedge m \dot{(O + x)} \nonumber \\
           &=& (Y - y + x) \wedge m \dot{(Y - y + x)} \nonumber \\
           &=& (Y + (x-y)) \wedge m \dot{(Y + (x - y))} = P_Y(X)
\end{subeqnarray}
where the last equation obviously is the same as eqn.~(\ref{momentum_written_with_O}), but
here Y takes the role of $O$ and $(x-y)$ takes the role of the vector $x$ in the old
description. So the same physical situation was described from another viewpoint. However,
the algebraic object $P$ is unchanged, i.e. $P_O(X) = P_Y(X)$. This clearly shows that $P$
does not depend on the chosen origin and that eqn.~(\ref{momentum_written_with_O}) is
indeed a coordinate free definition. \\

\emph{Second approach}: This one is not so explicit on being independent of the origin. We
just express $X$ by $Y$ using eqn.~(\ref{momentum_written_with_O}) and
(\ref{eq:identity_transformation_X_Y}) directly to write
\begin{subeqnarray}
    P_O(X) &=& X \wedge m \dot{X} \nonumber \\
           &=& (Y + (X-Y)) \wedge m \dot{(Y + (X-Y))} \nonumber \\
           &=& (Y + (x-y)) \wedge m \dot{(Y + (x - y))} = P_Y(X)
\end{subeqnarray}
We can again see that this is structurally identical to
eqn.~(\ref{momentum_written_with_O}) just seen from $Y$ instead from $O$. \\

The total momentum of a set of mass points $m_i X_i$ is defined as the sum of all momentum
blades.
\begin{equation}
    P_{total} = \sum_i{P_i}
    \label{eq:total_sum_momentum}
\end{equation}
$P_{total}$ will always be a momentum blade in $2D$ $PGA$, while in $3D$ $PGA$ (a
four-dimensional algebra) it could also result in a general bivector that might not be
factorizable into a blade (see \cite{Browne_Grassmann-Algebra_Vol1:2012}). That would be
the case, if $P \wedge P \ne 0 \; \ps$, i.e. when the result is a pseudoscalar that is not
equal to a zero multiple of the pseudoscalar \bv{1234}. This can always be used to
determine whether a bivector is a blade in $3D$ $PGA$. \\

Using points $X_i$ moved by a motor $M(t) = exp_\veedot(\frac{1}{2}B_w(t)) \veedot M_0$ in
the world frame as in eqn.~(\ref{eq:motor_pga_world_frame}) results in a local speed of
$\dot{X}_w = rcmt(\Omega_w,X_w)$ as shown in eqn.~(\ref{eq:rcmt_pga_world}). Therefore we
can write the total momentum as
\begin{eqnarray}
    P_{total} &=& \sum_i{P_i} \nonumber \\
              &=& \sum_i X_i \wedge m_i \dot{X}_i \nonumber \\
              &=& \sum_i X_i \wedge m_i \; rcmt(\Omega_w,X_i) \nonumber \\
    P_{total} &\equiv& I_w[\Omega_w]
              \label{eq:intertia_map_world_frame}
\end{eqnarray}
where eqn.~(\ref{eq:intertia_map_world_frame}) is a linear function of $\Omega_w$ that
maps the bivector rate in the world frame $\Omega_w$ to the total momentum $P_{total}$. It
is a function that takes a bivector as an argument and returns a bivector (``intertia
map'' $I_w$). It describes the influence of the mass distribution of the points on the
movement via the combined Newton/Euler law
\begin{equation}
    \dot{P}_{total}=F_{total}
    \label{eq:Newton_Euler_law}
\end{equation}
(for $F_{total}$ see eqn.~(\ref{eq:total_sum_forque})). If we go from discrete mass points
to a continuous mass distribution the sum has to be replaced by an integral over $dm =
\rho dV$. \\

If we want to switch to the body frame, we have to transform the components of
eqn.~(\ref{eq:intertia_map_world_frame}) accordingly using
eqn.~(\ref{eq:world_to_body_pga}), i.e. $\Omega_b = \rrev{M} \veedot \Omega_w \veedot M$
and eqn.~(\ref{eq:body_to_world_pga}), i.e. $\Omega_w  = M \veedot \Omega_b \veedot
\rrev{M}$ resulting in
\begin{equation}
    I_b[\Omega_b] \equiv \rrev{M} \veedot I_w[ M \veedot \Omega_b \veedot \rrev{M} ] \veedot M
    \label{eq:intertia_map_body_frame}
\end{equation}
The mapping to the classical intertia can be done when splitting $\Omega_b$ relative to
the centroid of the body with a projective point split on $\Omega_b$ and on
$I_b[\Omega_b]$ as shown by \cite{Dynamics_in_PGA_plane-based_Dorst-DeKennik:2023} on page
25f and adjusted to our approach using regressive geometric products.

Now we introduce some typically used reference points to describe our kinemetics (via the
bivector rate of change $\Omega$) and our dynamics (via the total momentum $P$):
\begin{itemize}
    \item \emph{Fixed body refence point} $O_b$ used as origin of the body frame. When
    the body is moved by the motor $M$ it will move in the world frame to $O_{b,w} = M
    \veedot O_b \veedot \rrev{M}$.
    \item \emph{Relative locations $r_i$ versus $O_b$} are relative positions in the body
    frame. Their projective counterpart in the body frame is $R_i = O_b + r_i$. When
    transformed to the world frame we get $X_i$ in the world frame by the usual
    transformation $X_i = M \veedot R_i \veedot \rrev{M} = M \veedot (O_b + r_i) \veedot
    \rrev{M}$.
    \item \emph{Center of mass} in the body frame is chosen to be located at $O_b$, such
    that $\sum_i{m_i r_i} = \vec{0}$. Thus the origin of the body frame is located at the
    center of mass. The center of mass in the world frame $C_w$ then is $C_w \equiv
    \sum_i{m_i X_i} / \sum_i{m_i} = \sum_i{m_i (M \veedot (O_b + r_i) \veedot \rrev{M})} /
    \sum_i{m_i} = M \veedot O_b \veedot \rrev{M}$
\end{itemize}


\newpage